\par\noindent\textbf{(1)} 在固定T,p下G总微分dG=Σ\_i μ\_i dn\_i，反应进度dξ令dn\_i=ν\_i dξ，故dG/dξ=Δ\_rG=Σ\_iν\_iμ\_i。平衡内部可逆反应要求Δ\_rG=0，且稳定平衡需沿反应坐标G的二阶曲率为正；不能只令各个μ\_i单独相等。
\par\noindent\textbf{(2)} 代入理想气体μ\_i=μ\_i°+RT ln(pi/p°)，得到Δ\_rG=Δ\_rG°+RT ln Qp，其中Qp=Π\_i(pi/p°)\textasciicircum{}\{ν\_i\}为无量纲反应商。平衡Qp=Kp，因此lnKp=−Δ\_rG°/(RT)，无需以量纲压力直接取对数。
\par\noindent\textbf{(3)} Gibbs–Helmholtz关系为∂(Δ\_rG°/T)/∂T=−Δ\_rH°/T²，因lnKp=−(Δ\_rG°/T)/R，得到d lnKp/dT=Δ\_rH°/(RT²)。只要求标准反应焓定义一致，若考虑较宽温区应计入ΔC\_p°使ΔH°随T变化。
\par\noindent\textbf{(4)} 若Δ\_rH°近常数，积分给ln[Kp(T2)/Kp(T1)]=(Δ\_rH°/R)(1/T1−1/T2)。对吸热Δ\_rH°>0且T2>T1右端为正，所以升温有利产物方向；放热反应反之，满足Le Chatelier原理。
